\BourbakiExercises{习题}

\begin{enumerate}
\def\labelenumi{\arabic{enumi})}
\tightlist
\item
  设 A 是一个不是体的主理想整环，P 是由 A 的不可约元素组成的一个代表系。A-模 M 是阿廷的，当且仅当它是挠模且存在 P 的一个有限子集 S，使得下列两个条件被满足：
\end{enumerate}

\begin{enumerate}
\def\labelenumi{(\roman{enumi})}
\item
  对每个 \(\pi \in \mathrm{P} - \mathrm{S}\) ，\(\pi\) -准素分量 \(\mathrm{M}_{\pi}\) （即 \(\mathrm{M}\) 的该分量，见 VII, §2, No.~2, p.~8）均归结为 0。
\item
  对每个 \(\pi\in S\) ，被 \(\pi\) 零化的 M 中元素组成的集合，视为体 A/ \(\pi\) A 上的向量空间，是有限维的。
\end{enumerate}

\begin{enumerate}
\def\labelenumi{\arabic{enumi})}
\setcounter{enumi}{1}
\tightlist
\item
  设 A 是环，M 是长度为 n 的有限长度 A-模，R 是 M 的一族对复合封闭的幂零自同态。
\end{enumerate}

\begin{enumerate}
\def\labelenumi{\alph{enumi})}
\item
  若 n = 1，则 R 的每个元素均为零。若 n \textgreater{} 1，证明存在 M 的一个子模 N，它异于 0 与 M，且在 R 的所有元素下稳定。（可以假设 R 有一个元素 \(s \neq 0\) ；选取一个使 \(s(M)\) 具有最小长度的，并证明 srs = 0 对每个 \(r \in R\) 成立。由此推出：可取 \(N = s(M)\) （若 \(Rs = \{0\}\) ），或取 \(N = \sum_{r \in R} r(s(M))\) （若 \(Rs \neq \{0\}\) 。）\(^{(1)}\)
\item
  由 a) 推出，存在 M 的一个若尔当--赫尔德列 \((\mathrm{M}_{i})_{0\leqslant i\leqslant n}\) ，使得 \(r(\mathrm{M}_{i})\subset\mathrm{M}_{i+1}\) 对所有 \(r\in R\) 及 \(0\leqslant i\leqslant n-1\) 成立。再推出：\(r_{1}\cdots r_{n}=0\) 对 R 的 n 个元素的每个序列 \((r_{i})_{1\leqslant i\leqslant n}\) 成立。
\end{enumerate}

\begin{enumerate}
\def\labelenumi{\arabic{enumi})}
\setcounter{enumi}{2}
\tightlist
\item
  \begin{enumerate}
  \def\labelenumii{\alph{enumii})}
  \tightlist
  \item
    给出一个交换域 K 的例子，以及一个同构 \(\varphi\) ，它从 K 映到 K 的某个子域 \(K'\) ，使得 K 在 \(K'\) 上是有限维的。在加法群 \(A = K \times K\) 上，我们通过令 \((x, y)(x', y') = (xx', xy' + y\varphi(x'))\) 定义一个环结构（II, §1, p.~382, 习题 7）。A 的左理想只有 \(\{(0, 0)\}\) 、 \(\{0\} \times K\) 以及 A，但对 K 的每个 \(K'\) -线性子空间 E，\(\{0\} \times E\) 都是 A 的右理想。由此推出 A 是左阿廷且左诺特的，但既非右阿廷也非右诺特。b) 给出一个有限长度模的例子，其自同态环既非左阿廷也非左诺特。
  \end{enumerate}
\end{enumerate}

\begin{enumerate}
\def\labelenumi{\alph{enumi})}
\setcounter{enumi}{2}
\tightlist
\item
  利用与 a) 相同的方法，给出一个左、右皆阿廷的环的例子，使其左长度异于右长度。
\end{enumerate}

\begin{enumerate}
\def\labelenumi{\arabic{enumi})}
\setcounter{enumi}{3}
\tightlist
\item
  设 \(\rho : \mathrm{A} \to \mathrm{B}\) 是一个环同态。
\end{enumerate}

\begin{enumerate}
\def\labelenumi{\alph{enumi})}
\item
  如果 B 作为右 A-模是自由的且非零，且环 B 是左阿廷的（相应地，左诺特的），那么环 A 是左阿廷的（相应地，左诺特的）。
\item
  给出一个例子，其中 A 是交换域，B 作为 A 上的右向量空间是有限维的，而 B 既非左阿廷也非左诺特（利用习题 3）。
\end{enumerate}

\begin{enumerate}
\def\labelenumi{\arabic{enumi})}
\setcounter{enumi}{4}
\item
  设 n 是一个整数 \(\geqslant 1\) 。环 A 是左阿廷的（相应地，左诺特的），当且仅当矩阵环 \(\mathbf{M}_{n}(\mathrm{A})\) 是左阿廷的（相应地，左诺特的）。若 A 是左阿廷的，则 \(\mathbf{M}_{n}(\mathrm{A})\) 的左长度等于 A 的左长度的 n 倍。
\item
  设 \(A\) 是环，\(e\) 是 \(A\) 中的一个幂等元。
\end{enumerate}

\begin{enumerate}
\def\labelenumi{\alph{enumi})}
\item
  设 \(\mathfrak{a}_1\) 是环 \(e\mathrm{A}e\) 的一个左理想，\(\mathfrak{a}\) 是 A 的由 \(\mathfrak{a}_1\) 生成的左理想。我们有 \(\mathfrak{a}_1 = e\mathfrak{a}e = \mathfrak{a} \cap e\mathrm{A}e\) 。
\item
  对 A 的每个左理想 a，有 \(a \cap eAe \subset eae\) ；给出一个使这是严格包含的例子（取 A 为矩阵环）。
\item
  对 A 的每个双边理想 a，有 \(a \cap eAe = eae\) 。
\item
  若环 A 是左阿廷的（相应地，左诺特的），则环 eAe 与 \((1-e)\mathrm{A}(1-e)\) 亦然。
\item
  给出一个例子，其中 eAe 与 \((1-e)\mathrm{A}(1-e)\) 都是交换域，而 A 容许双边理想的无穷严格递增序列与无穷严格递降序列。（取一个交换域 K 上的无限维向量空间 V，给加法群 \(A=K\times V\times K\) 赋予由 \((\lambda,x,\mu)(\lambda',x',\mu')=(\lambda\lambda',\lambda x'+\mu'x,\mu\mu')\) 定义的环结构，并令 \(e=(1,0,0)\) 。）
\end{enumerate}

\begin{enumerate}
\def\labelenumi{\arabic{enumi})}
\setcounter{enumi}{6}
\tightlist
\item
  \begin{enumerate}
  \def\labelenumii{\alph{enumii})}
  \tightlist
  \item
    设 A 是一个环， \(\mathcal{I}\) 是 A 中形如 \(Ae\) 的左理想组成的集合，其中 \(e\) 是 A 中的一个幂等元。证明下列性质等价：
  \end{enumerate}
\end{enumerate}

\begin{enumerate}
\def\labelenumi{(\roman{enumi})}
\item
  \(\mathcal{I}\) 的每个非空子集都有极小元。
\item
  \(\mathcal{I}\) 的每个非空子集都有极大元。
\item
  不存在 \(\mathcal{I}\) 中非零理想组成的无限族，使其和为直和。
\item
  不存在 A 中非零幂等元组成的无限族 \((e_i)\) ，使得 \(e_i e_j = 0\) 对 \(i \neq j\) 成立。
\end{enumerate}

\begin{enumerate}
\def\labelenumi{\alph{enumi})}
\setcounter{enumi}{1}
\item
  给出一个环 A 的例子，使其中每个由单生成左理想组成的非空集都有极大元和极小元（从而 \(a\) ) 中的等价条件得到满足），尽管 A 既非左阿廷又非左诺特（利用习题 3）。
\item
  设 M 是一个模。若 M 的有限生成子模组成的每个非空集都有极大元，则 M 是诺特的。
\item
  给出一个非阿廷模的例子，使其中每个由有限生成子模组成的非空集都有极小元。
\end{enumerate}

\begin{enumerate}
\def\labelenumi{\arabic{enumi})}
\setcounter{enumi}{7}
\item
  设 A 是环，B 与 D 是 A 的两个子环。证明：若 D 是体且环 B 是左阿廷的，则 \(B \cap D\) 是体。
\item
  设 A 是非零交换环，I 是无限集。环 \(\mathrm{A}[(\mathrm{X}_{i})_{i\in\mathrm{I}}]\) 与 \(\mathrm{A}[[(\mathrm{X}_{i})_{i\in\mathrm{I}}]]\) 都既非阿廷也非诺特。
\item
  设 A 是环，a 与 b 是 A 的两个双边理想。若环 A/a 与 A/b 都是左阿廷的（相应地，左诺特的），则环 A/a∩b 亦然；与此相对，环 A/ab 不一定是左阿廷的（相应地，左诺特的）。
\item
  设 A、B 是环，M 是有限生成的左 A-模，P 是有限生成的右 A-模，N 是一个 (A, B)-双模。若右 B-模 N 是阿廷的（相应地，诺特的），则右 B-模 \(\operatorname{Hom}_{\mathrm{A}}(\mathrm{M}, \mathrm{N})\) 与 \(P \otimes_{A} N\) 亦然。
\item
  设 \(\mathrm{K}\) 是交换域。用 \((e_n)_{n\in \mathbf{N}}\) 表示向量空间 \(\mathrm{V} = \mathrm{K}^{(\mathbf{N})}\) 的典范基。我们通过令下式来定义自同态的序列 \((u_n)_{n\in \mathbf{N}}\) （它们都是 \(\mathrm{V}\) 的自同态）：
\end{enumerate}

\[
\begin{array}{l l} u _ {0} (e _ {0}) = 0 & \qquad u _ {0} (e _ {1}) = 0 \qquad u _ {0} (e _ {n + 1}) = e _ {n} \text { for } n \geqslant 1 \\ u _ {m} (e _ {0}) = e _ {m} & \qquad u _ {m} (e _ {n}) = 0 \text { for } m \geqslant 1 \text { and } n \geqslant 1. \end{array}
\]

设 \(\mathrm{A}\) 是 \(\operatorname{End}_{\mathrm{K}}(\mathrm{V})\) 的由序列 \((u_n)_{n \in \mathbf{N}}\) 生成的含幺 K-子代数。证明左 A-模 \(\mathrm{V}\) 是阿廷的、有限生成的（甚至是单生成的），但不是诺特的。

\begin{enumerate}
\def\labelenumi{\arabic{enumi})}
\setcounter{enumi}{12}
\tightlist
\item
  \begin{enumerate}
  \def\labelenumii{\alph{enumii})}
  \tightlist
  \item
    设 \(\rho: \mathrm{A} \to \mathrm{B}\) 是环同态。假设 A 是左诺特的，且存在 B 的一个有限子集 S，其元素彼此交换且与 \(\rho(\mathrm{A})\) 的元素交换，使得环 B 由 \(\rho(\mathrm{A})\) 与 S 生成。则 B 是左诺特的。
  \end{enumerate}
\end{enumerate}

\begin{enumerate}
\def\labelenumi{\alph{enumi})}
\setcounter{enumi}{1}
\tightlist
\item
  设 A 是非零交换环。给出一个结合含幺 A-代数的例子，它由有限多个元素生成，而作为环既非左诺特也非右诺特（考察自由结合代数 \(\operatorname{Libas}_{\mathrm{A}}(\mathrm{I})\) （III, §2, No.~7, p.~449, 定义 2），其中 I 是基数 \(\geqslant 2\) 的集合）。
\end{enumerate}

\begin{enumerate}
\def\labelenumi{\arabic{enumi})}
\setcounter{enumi}{13}
\tightlist
\item
  设 V 是体 K 上可数无限维的向量空间。V 的有限秩自同态组成的集合 I 是环 \(B = \text{End}_{K}(V)\) 的一个双边理想。证明环 \(A = B/I\) 的双边理想只有 \(\{0\}\) 与 A，但 A 既非左诺特也非右诺特。
\end{enumerate}

¶ 15) 设 A 是交换环，M 是有限生成的 A-模。

\begin{enumerate}
\def\labelenumi{\alph{enumi})}
\item
  假设对 A 的理想的每个递增序列 \((\mathfrak{a}_{n})_{n\in\mathbf{N}}\) ，序列 \((\mathfrak{a}_{n}\mathrm{M})_{n\in\mathbf{N}}\) 都是平稳的。证明 A-模 M 是诺特的。（化归到 A-模 M 忠实、且对 A 的每个理想 \(a\neq0\) ，A-模 M/aM 都是诺特的情形。此时存在 M 的一个子模 \(N_{0}\) ，它在使 M/N 忠实的子模 \(N\subset M\) 当中是极大的。证明 A-模 \(M/N_{0}\) 是诺特的，再证明环 A 是诺特的，并由此完成证明。）
\item
  *假设对 A 的理想的每个递降序列 \((\mathfrak{a}_n)_{n \in \mathbf{N}}\) ，序列 \((\mathfrak{a}_n\mathrm{M})_{n \in \mathbf{N}}\) 都是平稳的。证明 A-模 M 是阿廷的。（化归到 A-模 M 忠实的情形。证明此时 A 只有有限多个极大理想，并且借助 VIII, p.~173 命题 1 的证明可知 A 的根是幂零的。推出存在 A 的极大理想的一个有限序列 \((\mathfrak{m}_{1}, \mathfrak{m}_{2}, \cdots, \mathfrak{m}_{n})\) ，其乘积零化 M。对 \(n.)_{*}\) 作归纳来完成证明。）
\end{enumerate}

\begin{enumerate}
\def\labelenumi{\arabic{enumi})}
\setcounter{enumi}{15}
\tightlist
\item
  \begin{enumerate}
  \def\labelenumii{\alph{enumii})}
  \tightlist
  \item
    设 A 是交换环，\(\rho: A \to B\) 是单射环同态，它使 B 成为有限生成的左 A-模。假设对 A 的理想的每个递增（相应地，递降）序列 \((\mathfrak{a}_{n})_{n \in \mathbf{N}}\) ，序列 \((\mathfrak{a}_{n}B)_{n \in \mathbf{N}}\) 都是平稳的（若环 B 是右诺特的（相应地，右阿廷的），则即属此情形）。证明环 A 是诺特的（相应地，阿廷的）（利用习题 15）。
  \end{enumerate}
\end{enumerate}

\begin{enumerate}
\def\labelenumi{\alph{enumi})}
\setcounter{enumi}{1}
\tightlist
\item
  给出一个左、右皆阿廷且诺特的环 B 及 B 的一个子环 A 的例子，使得 B 作为左、右 A-模都是有限生成的，而环 A 既非左也非右诺特或阿廷（选取一个适当的整环 C，其分式域为 K，取 B 为矩阵环 \(\mathbf{M}_{3}(\mathrm{K})\) ，并取 A 为 B 的由矩阵 \((a_{ij})_{1\leqslant i,j\leqslant3}\) 组成的子环，这些矩阵满足 \(a_{ij}=0\) （对 \(1\leqslant i<j\leqslant3\) ）且 \(a_{22}\in C.\) ）
\end{enumerate}

\begin{enumerate}
\def\labelenumi{\arabic{enumi})}
\setcounter{enumi}{16}
\tightlist
\item
  设 E 是一个幺半群，e 是其单位元，S 是 E 的一个子集，\(S'\) 是 E 的由 S 生成的子幺半群。我们说 S 在 E 上容许左分式演算，如果下列两个（当 E 交换时是重言式的）条件被满足：
\end{enumerate}

\begin{enumerate}
\def\labelenumi{(\roman{enumi})}
\item
  对每个 \(a \in \mathrm{E}\) 与每个 \(s \in \mathrm{S}\) ，存在 \(b \in \mathrm{E}\) 与 \(t \in \mathrm{S}'\) ，使得 \(ta = bs\) 。
\item
  对每组的 \(a \in \mathrm{E}\) 、 \(b \in \mathrm{E}\) 与 \(s \in \mathrm{S}\) ，若 \(as = bs\) ，则存在 \(t \in \mathrm{S}'\) ，使得 \(ta = tb\) 。
\end{enumerate}

\(a)\) 假设存在一个同态 \(u\) ，它从 E 映到幺半群 \(\mathrm{E}'\) ，使得 \(u(\mathrm{S})\) 的每个元素在 \(\mathrm{E}'\) 中可逆，且对 E 中的每对 \(a, b\) ，若 \(u(a) = u(b)\) ，则存在 \(s \in \mathrm{S}'\) ，使得 \(sa = sb\) 。则 S 在 E 上容许左分式演算。

\begin{enumerate}
\def\labelenumi{\alph{enumi})}
\setcounter{enumi}{1}
\item
  假设 S 在 E 上容许左分式演算。在 \(E \times S'\) 中，我们用 \((a, s) \sim (b, t)\) 表示关系``存在 E 中的 c 与 d，使得 ca = db 且 \(cs = dt \in S'\) 。''这是一个等价关系。我们令 \(S^{-1}E = (E \times S') / \sim\) ，并用 \(\varepsilon : E \to S^{-1}E\) 表示把 \(a \in E\) 映到 \((a, e)\) 的类的映射。在 \(S^{-1}E\) 上存在唯一的幺半群结构，使得 \(\varepsilon\) 是保单位元的同态，使得 \(\varepsilon(s)\) 对每个 \(s \in S\) 可逆，且使得对 \((a, s) \in E \times S'\) ，\((a, s)\) 在 \(S^{-1}E\) 中的类等于 \(\varepsilon(s)^{-1}\varepsilon(a)\) 。我们说 \(S^{-1}E\) 是 E 的相伴于 S 的左分式幺半群。当幺半群 E 交换时，它与 I, §2, No.~4, p.~18 中定义的那个一致。
\item
  设 \(a, b\) 在 E 中。\(\varepsilon(a) = \varepsilon(b)\) 成立，当且仅当存在 \(s \in S'\) ，使得 \(sa = sb\) 。映射 \(\varepsilon\) 是单射的（相应地，双射的），当且仅当 S 的每个元素是右正则的（相应地，可逆的）。
\item
  设 f 是从 E 到幺半群 F 的一个保单位元同态，使得 \(f(\mathrm{S})\) 的每个元素在 F 中可逆。存在唯一的同态 \(\overline{f}: S^{-1}E \to F\) ，使得 \(f = \overline{f} \circ s\) 。若 f 是单射，则 \(\overline{f}\) 是单射。
\item
  我们说 S 在 E 上容许右分式演算，如果它在 E 的反幺半群 \(E^{o}\) 上容许左分式演算。在这种情形，定义 E 的相伴于 S 的右分式幺半群 \(ES^{-1}\) ，注意到 \((\mathrm{ES}^{-1})^{\circ} = \mathrm{S}^{-1}(\mathrm{E}^{\circ})\) ，并对右分式幺半群改写 b)。若 S 在 E 上容许左、右分式演算，则幺半群 \(S^{-1}E\) 与 \(ES^{-1}\) 典范同构。
\end{enumerate}

\begin{enumerate}
\def\labelenumi{\arabic{enumi})}
\setcounter{enumi}{17}
\tightlist
\item
  \begin{enumerate}
  \def\labelenumii{\alph{enumii})}
  \tightlist
  \item
    设 A 是环，S 是 A 的一个子集，S' 是 A 的包含 1 与 S 且对乘法封闭的最小子集。我们说 S 在 A 上容许左分式演算，如果下列两个（当 A 交换时是重言式的）条件被满足：
  \end{enumerate}
\end{enumerate}

\begin{enumerate}
\def\labelenumi{(\roman{enumi})}
\item
  对每个 \(a \in \mathrm{A}\) 与每个 \(s \in \mathrm{S}\) ，存在 \(b \in \mathrm{A}\) 与 \(t \in \mathrm{S}'\) ，使得 \(ta = bs\) 。
\item
  对每组的 \(a \in \mathrm{A}\) 与 \(s \in \mathrm{S}\) ，若 \(as = 0\) ，则存在 \(t \in \mathrm{S}'\) ，使得 \(ta = 0\) 。
\end{enumerate}

这等价于说 S 在给集合 A 只赋予乘法而得的幺半群上容许左分式演算（习题 17）。证明：此时在乘法幺半群 \(S^{-1}A\) 上存在唯一的加法，使得赋有该加法及其乘法的 \(S^{-1}A\) 是一个环，且典范映射 \(\varepsilon: A \to S^{-1}A\) （同前所引）是环同态。我们说 \(S^{-1}A\) 是 A 的相伴于 S 的左分式环。当 A 交换时，它与 I, §8, No.~12, p.~113 中定义的环一致。

\begin{enumerate}
\def\labelenumi{\alph{enumi})}
\setcounter{enumi}{1}
\item
  集合 S 在 A 上容许左分式演算，当且仅当存在一个环 B 与一个环同态 \(u: A \to B\) ，使得 \(u(S)\) 的每个元素在 B 中可逆，且对每个 \(a \in u^{-1}(0)\) ，存在 \(s \in S\) ，使得 sa = 0。
\item
  假设 S 在 A 上容许左分式演算。典范同态 \(\varepsilon : \mathrm{A} \to \mathrm{S}^{-1}\mathrm{A}\) 的核是这样的 \(a \in \mathrm{A}\) 组成的集合：存在 \(s \in \mathrm{S}\) ，使得 \(sa = 0\) 。
\end{enumerate}

设 \(f\) 是从 A 到环 B 的一个同态，使得 \(f(\mathrm{S})\) 的每个元素在 B 中可逆。存在唯一的同态 \(\overline{f} : \mathrm{S}^{-1}\mathrm{A} \to \mathrm{F}\) ，使得 \(f = \overline{f} \circ s\) 。若 \(f\) 是单射，则 \(\overline{f}\) 是单射。

\begin{enumerate}
\def\labelenumi{\alph{enumi})}
\setcounter{enumi}{3}
\item
  对右分式演算叙述与 a)、b)、c) 类似的命题；若 S 在 A 上容许右分式演算，我们把 A 的右分式环记作 AS \(^{-1}\) 。在此情形，我们有 \((\mathrm{AS}^{-1})^{\circ} = \mathrm{S}^{-1}(\mathrm{A}^{\circ})\) 。若 S 在 A 上容许左、右分式演算，则环 S \(^{-1}\) A 与 AS \(^{-1}\) 典范同构。
\item
  设 A 是一个无零因子的非零环。则 A 容许一个左分式域（I, §9, p.~177, 习题 15），当且仅当 A 的非零元素组成的集合 S 在 A 上容许左分式演算，且在此情形，\(S^{-1}A\) 是 A 的左分式域。
\end{enumerate}

\begin{enumerate}
\def\labelenumi{\arabic{enumi})}
\setcounter{enumi}{18}
\tightlist
\item
  设 A 是环，S 是 A 的一个在 A 上容许左分式演算的子集（习题 18），S' 是 A 的包含 1 与 S 且对乘法封闭的最小子集。我们用 \(S^{-1}A\) 表示 A 的相伴于 S 的左分式环，并用 \(\varepsilon:A\to S^{-1}A\) 表示典范同态。
\end{enumerate}

\begin{enumerate}
\def\labelenumi{\alph{enumi})}
\item
  设 M 是一个 A-模。我们把 \(S^{-1}A\) -模 \(S^{-1}A \otimes_{A} M\) 记作 \(S^{-1}M\) ，并称之为 M 的相伴于 S 的左分式模。证明 \(S^{-1}M\) 的每个元素都可以写成 \(\varepsilon(s)^{-1} \otimes x\) ，其中 \(x \in M\) 且 \(s \in S'\) ；并且对 M 中的 x, y 与 \(S'\) 中的 s, t，关系 \(\varepsilon(s)^{-1} \otimes x = \varepsilon(t)^{-1} \otimes y\) 等价于存在 A 的一对元素 \((c, d)\) ，使得 cx = dy 且 \(cs = dt \in S'\) 。
\item
  我们用 \(\varepsilon_{M}\) 表示典范 A-线性映射 \(x \mapsto 1 \otimes x\) （从 M 到 \(S^{-1}M\) ）。它的像生成 \(S^{-1}A\) -模 \(S^{-1}M\) ；它的核由这样的 \(x \in M\) 组成：存在 \(s \in S'\) ，使得 sx = 0。映射 \(\varepsilon_{M}\) 是单射的（相应地，双射的），当且仅当对每个 \(s \in S\) ，M 的位似 \(x \mapsto sx\) 是单射的（相应地，双射的）。
\end{enumerate}

\(c)\) \(\mathrm{S}^{-1}\mathrm{M} = 0\) 成立，当且仅当对每个 \(x\in \mathrm{M}\) ，存在 \(s\in \mathrm{S}\) ，使得 \(sx = 0\) 。

\begin{enumerate}
\def\labelenumi{\alph{enumi})}
\setcounter{enumi}{3}
\item
  设 N 是 M 的一个子模。典范 \(S^{-1}A\) -线性映射（从 \(S^{-1}N\) 到 \(S^{-1}M\) ）是单射。（*换言之，右 A-模 \(S^{-1}A\) 是平坦的。*）我们把 \(S^{-1}N\) 与它在 \(S^{-1}M\) 中此映射下的像等同。M 的子模 \(\bar{\varepsilon}_{\mathrm{M}}^{-1}(\mathrm{S}^{-1}\mathrm{N})\) 由这样的 \(x \in M\) 组成：存在 \(s \in S'\) ，使得 \(sx \in N\) 。我们说它是 N 对 S 的左饱和。
\item
  映射 \(N^{\prime} \mapsto \varepsilon_{M}^{-1}(N^{\prime})\) 是一个从 \(S^{-1}A\) -子模组成的有序集（这些子模取自 \(S^{-1}M\) ）到 M 的对 S 左饱和（即等于其饱和）的 A-子模组成的有序集上的同构。逆同构是 \(N \mapsto S^{-1}N\) 。f) 若 A-模 M 是诺特的，则 \(S^{-1}A\) -模 \(S^{-1}M\) 是诺特的。若 A-模 M 是阿廷的，则 \(S^{-1}A\) -模 \(S^{-1}M\) 是阿廷的，且映射 \(\varepsilon_{M}: M \to S^{-1}M\) 是满射。
\end{enumerate}

\(g\) ) 若环 A 是左诺特的（相应地，左阿廷的），则环 \(\mathrm{S}^{-1}\mathrm{A}\) 亦然。\(h\) ) 对右模的右分式模，叙述与本习题类似的命题。

\begin{enumerate}
\def\labelenumi{\arabic{enumi})}
\setcounter{enumi}{19}
\item
  证明一个无零因子的左诺特环 A 容许左分式域（习题 18）（若 a 与 s 是 A 的非零元素，利用理想序列 \((As + Asa + \cdots + Asa^{n})_{n \in \mathbf{N}}\) 是平稳的这一事实，证明 \(Aa \cap As\) 不归结为 0）。
\item
  设 A 是环，\(\sigma\) 是环 A 的一个自同态，\(d\) 是 A 的加法群的一个满足关系式 \(d(ab) = \sigma(a)d(b) + d(a)b\) 的幂零自同态。\(a)\) 证明集合 \(S = \{X\}\) 在环 \(B = A[X]_{\sigma,d}\) （VIII, p.~11）上容许左分式演算（习题 18）。若 \(\sigma\) 是单射，则从 B 到 \(S^{-1}B\) 的典范同态是单射，且我们把 B 与 \(S^{-1}B\) 的一个子环等同。若 \(\sigma\) 是双射，则 S 在 B 上还容许右分式演算，且 \(S^{-1}B\) 的每个元素都可以唯一地写成 \(\sum_{n\in\mathbf{Z}}a_nX^n\) ，其中 \((a_n)_{n\in\mathbf{Z}}\) 是 A 的一个有限支集的元素族。对 \(a \in \mathrm{A}\) ，把元素 \(\mathrm{X}^{-1}a\) ——它属于 \(\mathrm{S}^{-1}\mathrm{B}\) ——写成这种形式。
\end{enumerate}

\begin{enumerate}
\def\labelenumi{\alph{enumi})}
\setcounter{enumi}{1}
\item
  证明环 \(A[X]_{\sigma,d}\) 中的乘法由连续性扩张为集合 \(A[[X]]\) （系数在 A 中、关于 X 的形式幂级数的集合）上的一个乘法。赋有该乘法的加法群 \(A[[X]]\) 是一个环，我们把它记作 \(A[[X]]_{\sigma,d}\) （若 d=0，也简记为 \(A[[X]]_{\sigma}\) ）。给出一个例子，其中一个常数项等于 1 的形式幂级数 \(u\in A[[X]]\) 在环 \(A[[X]]_{\sigma,d}\) 中既无左逆也无右逆。
\item
  证明 \(A[[X]]_{\sigma}\) 的可逆元素恰是那些常数项在 A 中可逆的元素。
\item
  把环 \(A[X]_{\sigma,d}\) 换为环 \(A[[X]]_{\sigma,d}\) 之后，叙述并证明 a) 的类似命题。
\item
  设 A 是体，\(\sigma\) 是 A 的一个自同构。则 \(\mathrm{S}^{-1}\mathrm{A}[[\mathrm{X}]]_{\sigma}\) （其中 \(\mathrm{S} = \{\mathrm{X}\}\) ）是一个体，其每个元素都可以唯一地写成 \(\sum_{n\in \mathbf{Z}}a_n\mathrm{X}^n\) ，其中 \((a_{n})_{n\in \mathbf{Z}}\) 是 A 的一个元素族，其指标 \(< 0\) 的非零项只有有限多个。我们把该体记作 \(\mathrm{A}((\mathrm{X}))_{\sigma}\) 。
\end{enumerate}

\(f\) ) 设 A 是交换域。描述体 \(\mathrm{A}((\mathrm{X}))_{\sigma}\) 的中心。由此推出一个在其中心上为无限次数的体的例子。

\begin{enumerate}
\def\labelenumi{\arabic{enumi})}
\setcounter{enumi}{21}
\tightlist
\item
  设 G 是群，H 是 G 的一个正规子群，使得商群 G/H 同构于 Z。设 g 是 G 的一个元素，它在 G/H 中的典范像生成 G/H。设 C 是交换环，A 是 H 在 C 上的代数 C{[}H{]}（III, §2, No.~6, p.~446），\((e_{h})_{h\in\mathrm{H}}\) 是其典范基，\(\sigma\) 是 A 的把 \(e_{h}\) 映到 \(e_{ghg^{-1}}\) 的自同构。
\end{enumerate}

\begin{enumerate}
\def\labelenumi{\alph{enumi})}
\item
  证明存在唯一的同态 \(u\) ，它从环 \(\mathrm{B} = \mathrm{A}[\mathrm{X}]_{\sigma}\) 到环 \(\mathrm{C}[\mathrm{G}]\) ，延拓 \(\mathrm{A} = \mathrm{C}[\mathrm{H}]\) 到 \(\mathrm{C}[\mathrm{G}]\) 中的典范单射，且把 \(\mathrm{X}\) 映为元素 \(e_g\) ——后者属于 \(\mathrm{C}[\mathrm{G}]\) 的典范基。
\item
  令 \(S = \{X\}\) 。证明 u 扩张为从环 \(S^{-1}B\) （参看习题 21）到 C{[}G{]}的一个同构。由此推出：若环 C{[}H{]}是左诺特的，则环 C{[}G{]}亦然。
\end{enumerate}

\begin{enumerate}
\def\labelenumi{\arabic{enumi})}
\setcounter{enumi}{22}
\tightlist
\item
  设 A 是环，I 是集合，\(\symbf{\sigma} = (\sigma_{i})_{i \in \mathrm{I}}\) 是环 A 的一族自同态，\(\symbf{d} = (d_{i})_{i \in \mathrm{I}}\) 是 A 的加法群的一族自同态，它们满足关系式
\end{enumerate}

\[
\sigma_ {i} \circ \sigma_ {j} = \sigma_ {j} \circ \sigma_ {i} \quad \sigma_ {i} \circ d _ {j} = d _ {j} \circ \sigma_ {i} \quad d _ {i} \circ d _ {j} = d _ {j} \circ d _ {i}
\]

对 \(i\in \mathrm{I},j\in \mathrm{I},i\neq j\) 以及

\[
d _ {i} (a b) = \sigma_ {i} (a) d _ {i} (b) + d _ {i} (a) b
\]

对 \(i \in I\) ， \(a \in A\) ， \(b \in A\) 。

\begin{enumerate}
\def\labelenumi{\alph{enumi})}
\tightlist
\item
  证明存在唯一的环 B，其加法群是多项式的加法群 A{[}X{]}——这些多项式以变元族 \(\mathbf{X} = (\mathrm{X}_{i})_{i \in \mathrm{I}}\) 、系数在 A 中——且其乘法具有下列性质：
\end{enumerate}

\begin{enumerate}
\def\labelenumi{(\roman{enumi})}
\item
  B 的元素 \(\mathrm{X}_i\) （ \(i \in \mathrm{I}\) ）两两可交换。
\item
  对每个 \(a \in A\) 与每个多重指标 \(\boldsymbol{\nu} = (\nu_i)_{i \in I} \in \mathbf{N}^{(I)}\) ，a 与一个由 \(\nu_i\) 个等于 \(X_i\) （对每个 \(i \in I\) ）的元素组成的有限序列按此顺序在 B 中的乘积是多项式 \(aX^\nu\) 。
\item
  对每个 \(a \in \mathrm{A}\) 与每个 \(i \in \mathrm{I}\) ，在环 B 中有关系式 \(\mathrm{X}_i a = \sigma_i(a)\mathrm{X}_i + d_i(a)\) 。
\end{enumerate}

\begin{enumerate}
\def\labelenumi{\alph{enumi})}
\setcounter{enumi}{1}
\item
  环 B 记作 \(A[X]_{\sigma,d}\) 。它具有下列泛性质：给定环 C、环同态 \(f: A \to C\) 以及 C 的一族两两可交换的元素 \((x_i)_{i \in I}\) ，使得 \(x_i f(a) = f(\sigma_i(a)) x_i + f(d_i(a))\) 对每个 \(a \in A\) 与每个 \(i \in I\) 成立，则存在唯一的环同态 \(g: B \to C\) ，它延拓 f，且使得 \(g(X_i) = x_i\) 对 \(i \in I\) 成立。
\item
  设 J 是 I 的一个子集。把环 \(\mathrm{A}[(\mathrm{X}_{i})_{i\in\mathrm{J}}]_{(\sigma_{i})_{i\in\mathrm{J}},(d_{i})_{i\in\mathrm{J}}}\) 记作 \(B'\) 。对每个 \(i\in I-J\) ，存在唯一的自同态 \(\sigma_{i}'\) （环 \(B'\) 的）与唯一的自同态 \(d_{i}'\) （\(B'\) 的加法群的），使得
\end{enumerate}

\[
\begin{array}{l l} \sigma_ {i} ^ {\prime} (a) = \sigma_ {i} (a) & d _ {i} ^ {\prime} (a) = d _ {i} (a) \qquad \text {for a\ensuremath{\in} A}, \\ \sigma_ {i} ^ {\prime} (\mathrm{X} _ {j}) = \mathrm{X} _ {j} & d _ {i} ^ {\prime} (\mathrm{X} _ {j}) = 0 \qquad \text {for j\ensuremath{\in} J}, \\ d _ {i} ^ {\prime} (\mathrm{PQ}) = \sigma_ {i} ^ {\prime} (\mathrm{P}) d _ {i} ^ {\prime} (\mathrm{Q}) + d _ {i} ^ {\prime} (\mathrm{P}) \mathrm{Q} & \text {for P\ensuremath{\in} B , Q\ensuremath{\in} B}. \end{array}
\]

令 \(\sigma' = (\sigma_i')_{i \in \mathbf{I} - \mathbf{J}}\) ，\(d' = (d_i')_{i \in \mathbf{I} - \mathbf{J}}\) 。从 \(\mathrm{A}[(\mathrm{X}_i)_{i \in \mathbf{I}}]\) 到 \(\mathrm{A}[(\mathrm{X}_j)_{j \in \mathbf{J}}][(\mathrm{X}_i)_{i \in \mathbf{I} - \mathbf{J}}]\) 的典范双射是从环 \(\mathrm{B} = \mathrm{A}[(\mathrm{X}_i)_{i \in \mathbf{I}}]_{\sigma, d}\) 到环 \(\mathrm{B}'[(\mathrm{X}_i)_{i \in \mathbf{I} - \mathbf{J}}]_{\sigma', d'}\) 的一个同构。

\(d)\) 若环 A 是左诺特的，则集合 I 有限；且若对每个 \(i \in \mathrm{I}\) ，A 的自同态 \(\sigma_i\) 都是双射，则环 \(\mathrm{B} = \mathrm{A}[\mathbf{X}]_{\sigma,d}\) 是左诺特的。\(e)\) 假设对每个 \(i \in \mathrm{I}\) ，态射 \(\sigma_i\) 是 A 上的恒等映射。设 E 是 A 的加法群的自同态环。证明存在唯一的环同态 \(\psi\) ，它从 \(\mathrm{B} = \mathrm{A}[\mathbf{X}]_{\sigma,d}\) 到 E，使得 \(\psi(\mathrm{X}_i) = d_i\) 对 \(i \in \mathrm{I}\) 成立，且对所有 \(a \in \mathrm{A}\) ，态射 \(\psi(a)\) 是 A 的左位似 \(b \mapsto ab\) 。\(f)\) 保留 \(e)\) 的假设。再设 A 是系数在一个（不必交换的）环 C 中的多项式环 \(\mathrm{C}[(\mathrm{T}_i)_{i \in \mathrm{I}}]\) ，且对每个 \(i \in \mathrm{I}\) ，态射 \(d_i\) 是 A 的导子 \(\mathrm{P} \mapsto \frac{\partial\mathrm{P}}{\partial\mathrm{T}_i}\) 。证明：若 C 非零且在 \(\mathbf{Z}\) 上无挠，则同态 \(\psi: \mathrm{B} \to \mathrm{E}\) 是单射。证明：若 C 是特征 \(p > 0\) 的环（V, §1, No.~2, p.~2），则 \(\psi\) 的核是 B 的由元素 \(\mathrm{X}_i^p\) （ \(i \in \mathrm{I}\) ）生成的左理想。

¶ 24) 设 G 是群，e 是其单位元，A 是非零交换环。用 A{[}G{]}表示群 G 在 A 上的代数（III, §2, No.~6, p.~446）。

\begin{enumerate}
\def\labelenumi{\alph{enumi})}
\item
  设 I 是有限集，\((\mathrm{G}_{i})_{i\in\mathrm{I}}\) 是 G 的一族子群，\((g_{i})_{i\in\mathrm{I}}\) 是 G 的一族元素。证明：若 G 是族 \((g_{i}\mathrm{G}_{i})_{i\in\mathrm{I}}\) 的并，则诸子群 \(G_{i}\) 之一在 G 中有有限指数。
\item
  假设 G 中异于 \(\{e\}\) 的每个共轭类都是无限的，且环 A 是整环。证明：若 a 与 b 是环 A{[}G{]}的两个非零双边理想，则我们有 \(ab \neq 0\) 。（在 a 与 b 中选取元素 \(a = \sum a_g g\) 与 \(b = \sum b_g g\) ，满足 \(a_e \neq 0\) 与 \(b_e \neq 0\) 。利用 a) 证明：必要时把 a 换成 \(hah^{-1}\) （取适当的 \(h \in G\) ），可以假设 \(a_g b_{g^{-1}} = 0\) 对 G 中每个 \(g \neq e\) 成立。）*c) 证明环 A{[}G{]}是左（相应地，右）阿廷的，当且仅当环 A 是阿廷的且群 G 是有限的。（若 A{[}G{]}是左阿廷的，由 VIII, p.~6 的定理 1 推出群 G 有有限长度；为证明它有限，化归到 A 是体且 G 是单群的情形，然后用 b) 与 VIII, p.~173 的命题 1。）*
\end{enumerate}

\begin{enumerate}
\def\labelenumi{\arabic{enumi})}
\setcounter{enumi}{24}
\tightlist
\item
  沿用习题 24 的记号。
\end{enumerate}

\begin{enumerate}
\def\labelenumi{\alph{enumi})}
\item
  对 G 的任一子群 H，令 \(I_{H}\) 为典范满射 \(\mathbf{Z}[G] \to \mathbf{Z}^{(G/H)}\) 的核。则映射 \(H \mapsto I_{H}\) 是从 G 的子群组成的集合到 \(Z[G]\) 的左理想组成的集合的递增单射；若 H 是 G 的正规子群，则理想 \(I_{H}\) 是双边的。
\item
  假设环 A{[}G{]}是左诺特的。则它也是右诺特的。对 G 的每个子群 H，环 A{[}H{]}是左诺特的；特别地，环 A 是诺特的。若 H 在 G 中正规，则环 A{[}G/H{]}是左诺特的。G 的每个子群都容许一个有限生成系。
\item
  证明下列性质等价：
\end{enumerate}

\begin{enumerate}
\def\labelenumi{(\roman{enumi})}
\item
  群 G 容许一个合成列，其商是有限的或同构于 Z。
\item
  群 G 容许一个合成列 \((\mathrm{G}_i)_{0\leqslant i\leqslant n}\) ，使得 \(\mathrm{G}_1\) 在 G 中有有限指数，且 \(\mathrm{G}_i / \mathrm{G}_{i + 1}\) 同构于 \(\mathbf{Z}\) （对 \(0\leqslant i\leqslant n - 1\) ）。
\end{enumerate}

当这些性质成立且环 A 是诺特的时，环 A{[}G{]}是左诺特的（利用习题 22）。

\begin{enumerate}
\def\labelenumi{\arabic{enumi})}
\setcounter{enumi}{25}
\item
  设 K 是交换域，A 是多项式环 K{[}T{]}，\(\sigma\) 是环 A 的自同态 \(\mathrm{P(T)}\mapsto\mathrm{P(T^{2})}\) 。环 \(A[X]_{\sigma}\) 不是左诺特的，尽管环 A 是诺特的。
\item
  设 D 是体，\(\sigma\) 是 D 的一个自同构，\(d\) 是 D 的加法群的一个满足 VIII, p.~9 中关系式 (1) 的自同态，A 是环 \(\mathrm{D}[\mathrm{X}]_{\sigma,d}\) （VIII, p.~11）。给定 A 的一个非零元素 \(\mathrm{P} = \sum_{n}a_{n}\mathrm{X}^{n}\) ，最大的整数 \(n\) （使 \(a_{n}\neq 0\) ）称为 P 的次数，记作 \(\deg (\mathrm{P})\) 。我们令 \(\deg (0) = -\infty\) 。
\end{enumerate}

\begin{enumerate}
\def\labelenumi{\alph{enumi})}
\item
  设 F, G 是 A 的元素且 \(G \neq 0\) ；证明存在 A 的完全确定的元素 Q, R（相应地， \(Q', R'\) ），使得 \(F = QG + R\) 且 \(\deg(R) < \deg(G)\) （相应地， \(F = GQ' + R'\) 且 \(\deg(R') < \deg(G)\) ）。
\item
  由此推出 A 的每个左（相应地，右）理想都是单生成的。若 F, G, H, K 是 A 的满足 AF + AG = AH 与 AF ∩ AG = AK 的元素，则我们有 \(\deg(\mathrm{F}) + \deg(\mathrm{G}) = \deg(\mathrm{H}) + \deg(\mathrm{K})\) （注意到 D-向量空间 A/AF 的维数是 \(\deg(\mathrm{F})\) ）。
\item
  A 的理想 AF 是极大的，当且仅当 F 不可约，即 F 不是两个非常数多项式的乘积。
\end{enumerate}

\begin{enumerate}
\def\labelenumi{\arabic{enumi})}
\setcounter{enumi}{27}
\tightlist
\item
  设 B 是环，A 是 B 的一个子环，x 是 B 的一个元素，使得 \(A + Ax = A + xA\) ，且环 B 由 \(A \cup \{x\}\) 生成。
\end{enumerate}

\begin{enumerate}
\def\labelenumi{\alph{enumi})}
\item
  设 M 是诺特的左 A-模；证明 B-模 \(\mathrm{M}_{(\mathrm{B})} = \mathrm{B} \otimes_{\mathrm{A}} \mathrm{M}\) 是诺特的。特别地，若环 A 是左诺特的，则 B 亦然。
\item
  设 A-模 M 有有限长度 n；证明 \(\mathrm{M}_{(\mathrm{B})}\) 的每个 B-子模都容许一个基数 \(\leqslant n\) 的生成子集（改编 VIII, p.~6 的定理 1 的证明）。
\end{enumerate}

